\section{Source audit, aim, and contribution}
\label{sec:audit}

\subsection{The inspected snapshot}

The repository was read through its GitHub connector at commit
\begin{center}\small\ttfamily
052fbaebfc955294b4c02f6532ab7e7a3f760f7a.
\end{center}
The proof-packet index lists 54 packets grouped by mathematical area.
It explicitly says that placing them in the repository did not verify,
rerun or compile their contents \cite{repoIndex}. We therefore distinguish
an author's stated result, a mathematical argument read in a source file,
a compiled host Lean theorem, and an actual object-iPRA derivation.

The review covered the index, the three previous framework packets, the
ConciseIPRA interfaces needed for the new Lean files, and selected
application reports. It is a targeted source audit, not a verification of
all 54 packets. The machine-readable manifest records inspected paths,
the pinned revision, available Git blob identifiers, and SHA-256 hashes
of materialized source text. Exact audit notes are included separately.

Three framework packets already exist at this revision:
\begin{enumerate}
\item \texttt{ipra\_framework\_2026-11-15}: an earlier compact framework.
\item \texttt{ipra\_embedding\_framework}: finite-demand transfer,
approximation and isolation, analytic regular branches, stable windows,
the odd-degree hyperelliptic norm count, and finite-algebra certificates.
\item \texttt{ipra\_framework}: finite contracts, an expanded analytic
chart constructor, finite-window arguments, and integral power-gap
certificates with explicit producers.
\end{enumerate}
Their reusable ideas are retained and cited as prior work
\cite{repoSmallFramework,repoEmbedding,repoFramework}. In particular,
bounded Booleanization, certificate composition, scalar Hensel steps,
ordinary bounded descent, and metastability are not presented as newly
discovered results of this report.

\subsection{The obstacles that determine the design}

\begin{longtable}{P{.20\linewidth}P{.34\linewidth}P{.37\linewidth}}
\toprule
Source & Specific difficulty or already available result & Consequence here\\
\midrule\endhead
ConciseIPRA & The equivalence transports an existing arithmetic derivation;
it does not turn arbitrary ambient Lean truth into one. & Every formal
endpoint is stated using \texttt{Prf PRAAx}.\\[4pt]
Existing frameworks & Finite observation contracts and bounded loops are
already developed. & Add exact restriction and coherence interfaces with
new worked algebraic constructions.\\[4pt]
Puiseux & General exact branch selection is impossible uniformly from
arbitrary Cauchy names; exact algebraic and regular-chart results exist.
& Solve simultaneous regular formal equations from a supplied invertible
Jacobian, with exact coherent jets. Retain the singular obstruction.\\[4pt]
Lagrange inversion; Bell numbers & Many identities concern finite
coefficient lists even when written using infinite generating functions.
& Give a common formal lifting and integral Hurwitz-series interface;
use their known identities as illustrations.\\[4pt]
Siegel general route & The proper-subspace fiber bound is described via
functions, divisors, and a zero/pole degree theorem. & Construct functions,
an exact dimension count and a sharp zero-count certificate for weighted
plane curves using finite polynomials.\\[4pt]
Earlier embedding framework & The fiber repair is already proved for
\(Y^2=f(X)\) with odd \(\deg f\). & Extend to \(Y^m-cX^d\) with arbitrary
lower-weight mixed terms and \(\gcd(m,d)=1\).\\[4pt]
Mordell--Weil & Contracting descent with explicit fuel already exists; actual
quotient representatives remain an input. & Reuse bounded descent only
as infrastructure. Do not claim it constructs those representatives.\\[4pt]
Faltings; Siegel & A general cardinal bound, height bound and complete list
are different targets. & Supply a counting component and state the
unresolved cover and height obligations separately.\\[4pt]
Diophantine geometry & Modular obstructions and bounded searches are finite;
passing local tests does not establish global rational points. & Add
compatible lifting through modulus powers; assert no local-to-global
converse.\\[4pt]
Prime distribution & Field certificates and major analytic estimates are
separate inputs. & The finite algebra modules can support certificates;
they do not prove the missing analytic estimates.\\
\bottomrule
\end{longtable}

The application rows refer to the inspected Puiseux, Lagrange-inversion,
Bell-number, Mordell--Weil, Faltings, Diophantine-geometry and prime-distribution
packets \cite{repoPuiseux,repoLagrange,repoBell,repoMW,repoFaltings,repoGeometry,repoPrimes}.

The most concrete source obligation addressed below is the paragraph
``The geometric input to the fibre bound'' in the Siegel report's
\texttt{paper/05\_general\_route.tex} \cite{repoSiegel}. Its original
description requests finite functions, independence data and a zero-count
argument. The earlier hyperelliptic repair settles a restricted instance.
Our weighted norm theorem settles a broader, explicitly stated family
without assuming a developed divisor theory. It does not settle the
separate finite-tuple approximation theorem required by that route.

\subsection{What this report adds}

The central framework addition is \emph{coherent finite presentation}.
Instead of immediately constructing an infinite object, construct its
requested finite observation, prove its finite specification, and prove
that restriction preserves that specification. A finite uniqueness theorem
then supplies coherence. This allows exact algebra, modular arithmetic and
formal differential algebra to use the same proof architecture as the
repository's metric approximation modules, while preserving their different
meanings.

The first main application is a filtered-algebra lifting theorem. For a
polynomial system, a root and an inverse Jacobian modulo the first
filtration level yield a fixed correction recurrence. Each step raises the
order of the residual by one; an explicitly nilpotent truncation therefore
needs only a prescribed finite number of steps. This gives formal implicit
functions in several unknowns and roots modulo powers of composite as well
as prime integer bases. A companion Hurwitz construction removes factorial
division from polynomial differential equations.

The second main application is an exact norm theorem for weighted plane
curves. It proves the determinant's degree and leading coefficient, then
shows that a fiber containing \(r\) distinct common zeros contributes
\(r\) factors to that determinant. The result gives a sharp point count,
an elementary dimension formula for an explicit space of functions, and a
proper-subspace fiber bound. An adjugate argument also gives an exact
function-field presentation for this curve class.

These use established algebraic ideas. Formal implicit-function theorems
and constructive Hensel methods have substantial prior literature
\cite{Hu,AlonsoLombardiNeuwirth}; norms as multiplication determinants are
standard \cite{StacksNorm}. The contribution claimed here is the explicit
generalization relative to the inspected packets, the integrated finite
proof architecture and its arithmetic audit. No priority claim over the
mathematical literature is made.

\subsection{Reading and proof-status conventions}

Sections~\ref{sec:foundations} and \ref{sec:refinement} give the common
rules. The next two chapters contain the main mathematical applications.
Section~\ref{sec:formalization} maps them to concrete Lean definitions,
compiled lemmas and remaining object derivations. The final chapter lists
exact accomplishments and boundaries.

The phrase ``has an iPRA proof'' in a mathematical theorem means that the
following human argument is expressed by the stated PR constructions and
bounded induction invariants. It does not mean that a corresponding Lean
proof object has already been checked. A compiled host theorem verifies
its own Lean proposition. A compiled theorem whose conclusion is
\texttt{Prf PRAAx ...} verifies an actual object derivation or a transformation
of supplied object derivations. These statuses are recorded individually;
compiling an unproved proposition definition is never counted as proving it.
